\documentclass[12pt,a4paper]{article}

% Pakiety podstawowe
\usepackage[utf8]{inputenc}
\usepackage[T1]{fontenc}
\usepackage[polish]{babel}
\usepackage{geometry}
\usepackage{titling}
\usepackage{enumitem}
\usepackage{hyperref}
\usepackage{tabularx}

% Ustawienia marginesów
\geometry{
    a4paper,
    left=25mm,
    right=25mm,
    top=25mm,
    bottom=25mm,
}

% Ustawienia linków
\hypersetup{
    colorlinks=true,
    linkcolor=blue,
    filecolor=magenta,      
    urlcolor=cyan,
    pdftitle={Zadanie 1 - Analiza Problemu Decyzyjnego},
}

% --- DANE DO TYTUŁU ---
\title{\textbf{Zadanie 1: Analiza problemu decyzyjnego}\\
\large Przedmiot: Systemy Wspomagania Decyzji (SWD)\\
\vspace{0.5cm}
\Large Temat: Optymalizacja lokalizacji, rozmiaru i typu schronów oraz miejsc bezpiecznych dla aglomeracji miejskiej (1-2 mln mieszkańców)}

\author{
    1. [Filip Filipkowski, 327344] \\
    2. [Antoni Kingston, 327284] \\
    3. [Hubert Sobociński, 327315]
}
\date{}

\begin{document}

\maketitle

\section*{Wstęp - Od problemu do decyzji}
Projekt zakłada opracowanie koncepcji Systemu Wspomagania Decyzji (SWD) wspierającego władze dużego miasta w strategicznym planowaniu infrastruktury ochronnej na wypadek wystąpienia zagrożeń (np. konflikt zbrojny, katastrofy przemysłowe). Poniżej zestawiono odpowiedzi na zasadnicze pytania kierunkowe:

\begin{itemize}
    \item \textbf{Kto podejmuje decyzję?} Decyzję podejmują władze samorządowe reprezentowane przez Prezydenta Miasta we współpracy z Wydziałem Zarządzania Kryzysowego i Bezpieczeństwa.
    \item \textbf{Jaki jest zakres wyboru?} Wybór polega na określeniu optymalnej konfiguracji sieci schronów (rozmieszczenie punktów na mapie miasta, wybór ich pojemności oraz określenie typu obiektu: adaptacja infrastruktury podziemnej czy budowa nowych jednostek).
    \item \textbf{Co oznacza dobra decyzja?} Dobra decyzja maksymalizuje procent populacji z dostępem do bezpiecznego schronienia w czasie pieszego dojścia do 15 minut, przy jednoczesnym zmieszczeniu się w budżecie inwestycyjnym.
    \item \textbf{Jakie są ograniczenia?} Ograniczenia wynikają z budżetu, uwarunkowań geologicznych, zagospodarowania przestrzennego oraz uwarunkowań prawnych.
\end{itemize}

\section{Szczegółowa klasyfikacja problemu decyzyjnego}

Zgodnie z metodyką systemów wspomagania decyzji, analizowany problem został zdekomponowany na następujące elementy składowe:

\subsection{Decydent}
Wydział Zarządzania Kryzysowego Urzędu Miasta, w szczególności Zespół ds. Planowania Przestrzennego i Infrastruktury Krytycznej (podejmowanie decyzji o charakterze strategicznym i wielomilionowych inwestycjach).

\subsection{Interesariusze}
Oprócz decydenta na kształt i odbiór decyzji wpływają następujące grupy:
\begin{itemize}
    \item \textbf{Mieszkańcy:} Główni beneficjenci systemu, zainteresowani krótkim czasem dojścia do schronu oraz brakiem uciążliwości budowy w najbliższym otoczeniu.
    \item \textbf{Państwowa Straż Pożarna i Obrona Cywilna:} Odpowiadają za ewidencję budowli ochronnych i organizację ewakuacji, dostarczają wytycznych technicznych.
    \item \textbf{Wojewoda i MSWiA:} Nadzorują realizację programu ochrony ludności i współfinansują inwestycje z budżetu państwa.
    \item \textbf{Zarządcy obiektów:} Operator metra, właściciele parkingów podziemnych i garaży, których obiekty mogą zostać zaadaptowane na cele ochronne.
    \item \textbf{Rady dzielnic:} Reprezentują lokalne społeczności, mogą wspierać lub blokować konkretne lokalizacje (syndrom NIMBY).
\end{itemize}

\subsection{Cel decyzji}
Optymalizacja rozmieszczenia infrastruktury ochronnej polegająca na:
\begin{itemize}
    \item \textbf{Maksymalizacji pokrycia populacji:} Zapewnieniu dostępu do schronu jak największej liczbie obywateli.
    \item \textbf{Minimalizacji kosztów i czasu budowy:} Zrównoważeniu drogich nowych konstrukcji z tańszą adaptacją (np. stacje metra, parkingi podziemne).
\end{itemize}

\subsection{Warianty decyzyjne}
Warianty stanowią unikalne kombinacje (scenariusze inwestycyjne) przypisujące lokalizacjom przestrzennym konkretne parametry:
\begin{itemize}
    \item \textbf{Lokalizacja:} Przydział konkretnych działek w dzielnicach lub istniejących obiektów inżynieryjnych.
    \item \textbf{Rozmiar/Pojemność:} Zróżnicowanie na małe schrony (do 50 osób, szybkie dojście) oraz duże obiekty masowe (powyżej 1000 osób, trudniejsze dojście).
    \item \textbf{Typ:} Budowa od zera (tzw. bunkry dedykowane) vs. adaptacja podwójnego przeznaczenia (parkingi, podziemia, metro).
\end{itemize}

\subsection{Kryteria oceny}
Poszczególne warianty inwestycyjne oceniane będą według następujących kryteriów:
\begin{enumerate}
    \item \textbf{Czas dostępu (izochrony dojścia pieszo):} Główny priorytet, czas mierzony w oparciu o siatkę ulic (max. 10-15 minut).
    \item \textbf{Nakłady finansowe:} Koszt projektowania, budowy/adaptacji oraz późniejszego utrzymania rocznego.
    \item \textbf{Pojemność całkowita:} Stosunek ilości dostępnych miejsc do gęstości zaludnienia w danym sektorze miasta.
    \item \textbf{Kategoria ochrony:} Stopień ochrony przed konkretnymi typami zagrożeń (np. promieniowanie, fala uderzeniowa).
\end{enumerate}

\subsection{Ograniczenia}
\begin{itemize}
    \item \textbf{Finansowe:} Maksymalny limit budżetu na dany rok (np. dotacje z budżetu państwa).
    \item \textbf{Przestrzenne i geologiczne:} Zwarte centra miast uniemożliwiają budowę obiektów od zera, występowanie wód gruntowych blokuje użycie niektórych terenów.
    \item \textbf{Prawne:} Wymagania prawa budowlanego oraz wytyczne Obrony Cywilnej dotyczące wentylacji czy wyjść ewakuacyjnych.
\end{itemize}

\subsection{Otoczenie decyzji}
Problematyka jest głęboko osadzona w dynamice dużego miasta (1-2 mln ludności). Otoczenie charakteryzuje się:
\begin{itemize}
    \item \textbf{Dobowymi migracjami:} Populacja w centrum (biurowce) i na osiedlach sypialniach drastycznie różni się w zależności od pory dnia i nocy.
    \item Złożonością miejskiej sieci komunikacyjnej (np. szerokie rzeki stanowiące barierę dla pieszych).
    \item Potencjalnym oporem społecznym (tzw. syndrom NIMBY - \textit{Not In My Back Yard}).
\end{itemize}

\subsection{Źródła danych}
System będzie zasilany danymi z następujących źródeł:
\begin{itemize}
    \item \textbf{GUS (NSP 2021):} Liczba ludności w siatce kilometrowej i podział na rejony statystyczne.
    \item \textbf{BDOT10k i OpenStreetMap:} Sieć ulic i ciągów pieszych, obrys budynków, przeszkody terenowe (rzeki, tory).
    \item \textbf{Dokumenty planistyczne (MPZP, Studium):} Przeznaczenie terenów i dostępne działki.
    \item \textbf{Mapy hydrogeologiczne (PIG-PIB):} Poziom wód gruntowych i warunki gruntowe.
    \item \textbf{Dane operatorów komórkowych i transportu publicznego:} Rozkład populacji w ciągu doby (scenariusz dzień/noc).
    \item \textbf{Rejestr budowli ochronnych PSP:} Istniejące schrony i ukrycia wraz z ich stanem technicznym.
\end{itemize}

\subsection{Niepewność}
Z uwagi na naturę problemu, system musi radzić sobie z ryzykiem i brakiem pełnej informacji:
\begin{itemize}
    \item Niepewność pory wystąpienia ewentualnego zdarzenia kryzysowego (determinuje, gdzie znajduje się większość obywateli miasta).
    \item Tempo migracji ludności w latach objętych scenariuszami inwestycyjnymi (zmiany demograficzne w dzielnicach).
    \item Nieprzewidywalność rzeczywistego uszkodzenia szlaków komunikacyjnych powstrzymujących obywateli przed dotarciem do schronu.
\end{itemize}

\subsection{Klasyfikacja problemu}
Opierając się na omawianych na zajęciach współczesnych metodach klasyfikacji SWD:
\begin{itemize}
    \item \textbf{Typ problemu:} Problem złożony o charakterze lokalizacyjnym (Location-Allocation Problem) oraz wielokryterialnym.
    \item \textbf{Podejście systemu:} System będzie miał charakter \textbf{hybrydowy}. 
    \item Połączy podejście \textbf{Data-driven / GIS} (wykorzystanie danych przestrzennych dotyczących demografii i infrastruktury) z podejściem \textbf{Model-driven} (algorytmy optymalizacyjne do wyznaczania tras i przypisania populacji, wsparte metodą AHP do priorytetyzacji dzielnic i kryteriów decyzji).
\end{itemize}

\section{Ogólna architektura systemu}
Proponowany system wpisuje się w klasyczną strukturę SWD, składającą się z trzech współpracujących podsystemów:
\begin{itemize}
    \item \textbf{Podsystem danych:} Baza danych przestrzennych (GIS) integrująca dane demograficzne, sieć komunikacyjną, ograniczenia terenowe oraz rejestr istniejących obiektów.
    \item \textbf{Podsystem modeli:} Algorytmy optymalizacji lokalizacji wyznaczające rozmieszczenie schronów i przypisanie do nich populacji, uzupełnione metodą AHP do ustalania wag kryteriów.
    \item \textbf{Interfejs użytkownika:} Mapa prezentująca warianty rozmieszczenia, porównanie scenariuszy (np. dzień/noc) oraz analiza typu \textit{what-if} (np. zmiana budżetu, wyłączenie mostu).
\end{itemize}

\end{document}
